% Part:sets-functions-relations
% Chapter: sets
% Section: schroder-bernstein

\documentclass[../../../include/open-logic-section]{subfiles}

\begin{document}

\olfileid{sfr}{siz}{sb}

\olsection{سائز دا تصور تے شروڈر-برنسٹین}

\begin{explain}
ایہ اک بدیہی خیال اے: جے $A$،~$B$ توں وڈا نہیں تے $B$،~$A$ توں
وڈا نہیں، تاں $A$ تے~$B$ ہم تعداد نیں۔ سچ پچھو تاں جے ایہ خیال
\emph{غلط} ہوندا، تاں اسی شاید ایہ گل جائز ای نہ ٹھہرا سکدے کہ ہم
تعدادیت دا ساڈا تعریف کیتا تصور سیٹاں دے ”سائزاں“ دے موازنے نال کوئی
واسطہ رکھدا اے!{} پر خوش قسمتی نال ایہ بدیہی خیال درست اے۔ ایس دا
جواز شروڈر-برنسٹین دا قضیہ دیندا اے۔
\end{explain}

\begin{thm}[شروڈر-برنسٹین]
	\ollabel{thm:schroder-bernstein}
	جے $\cardle{A}{B}$ تے $\cardle{B}{A}$،
	تاں $\cardeq{A}{B}$۔
\end{thm}

\begin{explain}
ہور لفظاں وچ، جے $A$ توں~$B$ ول !!a{injection}، تے
$B$ توں~$A$ ول !!a{injection} موجود اے، تاں $A$ توں~$B$ ول
!!a{bijection} موجود اے۔

پر ایس نتیجے نوں ثابت کرنا واقعی خاصا \emph{اوکھا} اے۔ اصل وچ کانتور
نے نتیجہ بیان تاں کیتا، پر اوہنوں ہورناں نے ثابت کیتا۔\footnote{
تاریخ بارے ہور جانن لئی مثلاً \citet[pp.~165--6]{Potter2004} ویکھو۔}
\oliflabeldef{sfr:cardinals:card-sb:sec}{اسی شروڈر-برنسٹین نوں
\olref[sfr][cardinals][card-sb]{sec} وچ جا کے ای \emph{ثابت} کرن دی
حالت وچ ہواں گے۔}{}%
فی الحال تسی ایس نوں بھروسے اُتے منّ سکدے او (تے منّنا پئے گا)۔

خوش قسمتی نال شروڈر-برنسٹین \emph{درست} اے، تے ایہ ایس سوچ نوں
درست ٹھہراؤندا اے کہ ساڈے تعریف کیتے تعلق، یعنی
$\cardeq{A}{B}$ تے $\cardle{A}{B}$، ”سائز“ نال کوئی واسطہ رکھدے
نیں۔ ایس توں ودھ، شروڈر-برنسٹین بہت \emph{فائدہ مند} اے۔ دو ہم
تعداد سیٹاں دے وچکار !!a{bijection} سوچنا اوکھا ہو سکدا اے۔
شروڈر-برنسٹین دا قضیہ سانوں موازنے نوں صورتاں وچ ونڈن دیندا اے، ایس
لئی سانوں صرف پہلے سیٹ توں دوجے ول !!a{injection}، تے اوسے طرح
دوجے توں پہلے ول، سوچنا پیندا اے۔
\end{explain}
\end{document}
